% Coulomb's law: the force on q due to Q, and the vectors that locate them.
%
% Pulled in with \tikzfig{ch01-coulomb-law.tex}. A fragment, not a document --
% one tikzpicture, no preamble -- so it inherits the fonts of whatever \inputs
% it. The libraries it needs are loaded by the preamble, not here -- see the
% note above the picture.
%
% Everything is placed by TikZ's own machinery rather than by measured offsets,
% so the drawing survives being resized or re-lettered:
%
%   - the charges and the origin are nodes, so every arrow is clipped at their
%     border automatically; `outer sep' is the gap, one number per kind of
%     object rather than one per arrow;
%   - the vector names are nodes on their own path with `auto', so TikZ chooses
%     which side of the line to sit on and the node's own inner sep sets how
%     far off it sits;
%   - the charge names are `label's, which stay put relative to the ring if the
%     ring is ever resized.
%
% What is left to choose is the style of the figure as a whole -- its scale, how
% big a charge is, how heavy the strokes are -- which is the level worth having
% knobs at. The one computed coordinate is the tip of F, and that is the point:
% it is R extended past q, so it cannot drift out of parallel with R.
%
% Both charges are positive, which is what puts F along R rather than against
% it. The charges sit off-axis so no direction looks preferred, and Q sits above
% the line O-q, which opens the triangle O-Q-q wide enough to letter on all
% three sides.
% Needs the arrows.meta and calc libraries. They are NOT loaded
% here: \usetikzlibrary inside a figure environment sets its "loaded" flag
% globally but its definitions only for that group, which silently disarms the
% next figure's load. gwbook.cls and _wrapper.tex load them in the preamble
% instead -- add to those two if this figure ever needs another.
\begin{tikzpicture}[
    % The figsize knob: distances scale, lettering and node sizes do not.
    scale=1.2,
    >={Latex},
    charge/.style = {circle, draw, fill=white, line width=0.9pt,
                     minimum size=12pt, inner sep=0pt, outer sep=2.5pt,
                     font=\normalsize},
    point/.style  = {circle, fill, minimum size=4.5pt, inner sep=0pt,
                     outer sep=0pt},
    vec/.style    = {->, line width=0.9pt},
    % F is the one thing the law asserts; the rest of the figure constructs it.
    force/.style  = {vec, line width=1.4pt},
    % outer sep already holds everything off the ring -- anchors sit outside
    % it, so a label inherits that clearance and needs no distance of its own.
    every label/.style = {inner sep=0pt},
  ]

  \node[point,  label=below right:$O$] (O) at (0,0)       {};
  \node[charge, label=below right:$Q$] (Q) at (1.15,1.95) {$+$};
  \node[charge, label=below right:$q$] (q) at (4.60,1.30) {$+$};

  % `auto' puts a name to the left of the way the arrow travels; `swap' takes
  % the other side. Here that is the choice of outside the triangle, in each
  % case, and it needs no coordinates.
  \draw[vec] (O) -- node[auto]       {$\mathbf{r}'$} (Q);
  \draw[vec] (O) -- node[auto, swap] {$\mathbf{r}$}  (q);

  % Set along the arrow: the expression is too wide to sit level over a sloped
  % line without fouling it.
  \draw[vec] (Q) -- node[auto, sloped] {$\mathbf{R} = \mathbf{r} - \mathbf{r}'$} (q);

  % R continued past q by 45% of its length -- which is the whole content of
  % the law, and is why the direction is taken from R instead of typed out.
  \draw[force] (q) -- node[auto] {$\mathbf{F}$} ($(q)!-0.45!(Q)$);

\end{tikzpicture}
